\documentclass[../../main.tex]{subfiles}

\begin{document}

\section{Arrays and Vectors}

One of the most important and essential topic in computer science
in general could be arrays and vectors and they provide a way
to effectively store and manipulate data. Let's get started with
arrays!

\subsection{Arrays}

A simple array of $n$ elements can be thought of as an $n$-tuple.

\begin{definition}{}{}
An \textbf{array} is a collection of elements of the same
data type.
\end{definition}

When we declare an array, we need to specify the type, name, and
the number of elements \cite{CS02-4}.

\begin{basicbox}
\begin{lstlisting}[style=basic]
dataType arrayName[number];
\end{lstlisting}
\end{basicbox}

Consider the following example.

\begin{cppbox}{CS0201.26053-01}{CS0201.26053-01}
\begin{lstlisting}[style=cppstyle]
#include <iostream>

int main() {
    int arr[5] = {1, 2, 3, 4, 5};

    for (int i = 0; i < 5; i++) {
        std::cout << arr[i] << " ";
    }
    std::cout << "\n";

    return 0;
}
\end{lstlisting}
\end{cppbox}

In the example above, we have an array \verb|arr| with five elements.
When we actually print them to the console with a for loop, we can
see that we have to start from \verb|i = 0| since the index of the
first element is $0$. To access an element, we can have the name
of the array, bracket, and the index number.

\begin{outputbox}
\begin{lstlisting}[style=basic]
1 2 3 4 5
\end{lstlisting}
\end{outputbox}

We can also edit an element by specifying it.

\begin{cppbox}{CS0201.26053-02}{CS0201.26053-02}
\begin{lstlisting}[style=cppstyle]
#include <iostream>

int main() {
    int arr[5] = {1, 2, 3, 4, 5};
    arr[2] = -3;
    std::cout << arr[2] << "\n";

    return 0;
}
\end{lstlisting}
\end{cppbox}

The line \verb|arr[2] = -3| will override the previous assignment.

\begin{outputbox}
\begin{lstlisting}[style=basic]
-3
\end{lstlisting}
\end{outputbox}

Continuing, an array need not be $1$-dimensional. There exists
multidimensional arrays that are made of lower dimensional
arrays \cite{CS02-6}. Below is an example of a $2 \times 3$
matrix.

\begin{cppbox}{CS0201.26053-03}{CS0201.26053-03}
\begin{lstlisting}[style=cppstyle]
#include <iostream>

int main() {
    int arr[2][3];

    for (int i = 0; i < 2; i++) {
        for (int j = 0; j < 3; j++) {
            arr[i][j] = (i + 1)*10 + (j + 1);
        }
    }

    for (int i = 0; i < 2; i++) {
        for (int j = 0; j < 3; j++) {
            std::cout << arr[i][j] << " ";
        }
        std::cout << "\n";
    }

    return 0;
}
\end{lstlisting}
\end{cppbox}

In the example above, we declared a multidimensional array with
two rows and three columns. Then, we used loops to assign a value
in each index. Using the same way to put an input, we printed the
values to the console. The index for the matrix should be familiar.

\begin{outputbox}
\begin{lstlisting}[style=basic]
11 12 13
21 22 23
\end{lstlisting}
\end{outputbox}

From the examples above, you may have noticed that array is not
really flexible in terms of number of elements. Sometimes we
might want to add or remove elements from an array, but it won't
allow that since the size is fixed. This is were vectors come
to place.

\subsection{Vectors}

A simple way of thinking vectors is resizable array.

\begin{definition}{}{}
A \textbf{vector} is a dynamic container that contains elements of
same the data type.
\end{definition}

Vectors are provided by the Standard Template Library and we need
to include directive to access them \cite{CS02-5}. Because
vectors are dynamic, we only need the data type and name of the
vector when declaring.

\begin{outputbox}
\begin{lstlisting}[style=basic]
std::vector<dataType> vectorName;
\end{lstlisting}
\end{outputbox}

Below is a quick example of declaring a vector and two ways of
printing the elements in the console.

\begin{cppbox}{CS0201.26053-0}{CS0201.26053-0}
\begin{lstlisting}[style=cppstyle]
#include <iostream>
#include <vector>

int main() {
    std::vector<int> vec = {1, 2, 3, 4, 5};

    for (int i = 0; i < 5; i++) {
        std::cout << vec[i] << " ";
    }
    std::cout << "\n";

    for (int element : vec) {
        std::cout << element << " ";
    }
    std::cout << "\n";

    return 0;
}
\end{lstlisting}
\end{cppbox}

From the example above, we declared a vector \verb|vec| and printed
the elements to the console in two different ways. Both methods are
fine and the \verb|:| symbol can be interpreted as ``in.''

\begin{outputbox}
\begin{lstlisting}[style=basic]
1 2 3 4 5
1 2 3 4 5
\end{lstlisting}
\end{outputbox}

With the idea in mind, let's discuss few key functions that
distinguishes a vector and an array.

The first function is \verb|.push_back()| function. This function
allows us to add an element as the last element. If you wish to
delete the last element, we can use the \verb|.pop_back()| function.
Lastly, \verb|.size()| will allow us to see the size of a vector.
Consider the following example.

\begin{cppbox}{CS0201.26053-05}{CS0201.26053-05}
\begin{lstlisting}[style=cppstyle]
#include <iostream>
#include <vector>

int main() {
    std::vector<int> vec = {1, 2, 3, 4, 5};
    
    vec.push_back(6);
    for (int x : vec) {
        std::cout << x << " ";
    }
    std::cout << "\n";
    std::cout << vec.size() << "\n";

    vec.pop_back();
    for (int x : vec) {
        std::cout << x << " ";
    }
    std::cout << "\n";
    std::cout << vec.size() << "\n";

    return 0;
}
\end{lstlisting}
\end{cppbox}

In this example, we declare a vector, add $6$ to the end, list the
elements, and find the size. Then, we will revert it to the original
vector by removing the last element that was added.

\begin{outputbox}
\begin{lstlisting}[style=basic]
1 2 3 4 5 6
6
1 2 3 4 5
5
\end{lstlisting}
\end{outputbox}

This is it for the C++ basics! There are so much more topics to
discuss in C++, but for now, let's solve some practice problems
before discussing more fundamental concepts in DSA.

\end{document}


